% Source: https://artofproblemsolving.com/community/c3h1597230_questions
% Stable identity: aops-topic-1597230-q02
\textbf{The expansion.} $\dfrac{1}{75}=\dfrac{4}{300}=\dfrac{1.333\ldots}{100}=0.01333\ldots$: after the digits $0$ and $1$, every digit is $3$. Indeed $0.01\overline{3}=\dfrac{1}{100}+\dfrac{1}{300}=\dfrac{4}{300}=\dfrac{1}{75}$.

So the $100$th digit after the decimal point is $\boxed{3}$.

\textbf{General method.} Write the decimal as a non-repeating block of $p$ digits followed by a repeating block of length $\ell$ (here $p=2$, the block ``$01$'', and $\ell=1$, the block ``$3$''). For $n>p$, the $n$th digit is the $r$th digit of the repeating block, where $r$ is the remainder of $n-p$ on division by $\ell$ (using $r=\ell$ when the remainder is $0$). For example $\dfrac{1}{7}=0.\overline{142857}$ has $p=0$, $\ell=6$; the $100$th digit is the $4$th digit of $142857$ (as $100=6\cdot16+4$), namely $8$.
